\section{Оговорка и приглашение}
\label{sec:disclaimer}

Честности ради оговоримся: крипто кома \(\ccoma\) --- это \emph{конечное} натуральное
число, пусть и невообразимо большое. Называя его «мощностью» и одевая в одежды теории
множеств, мы играем на двойном смысле слова: и «кардинальное число», и «сила». В этом и
состоит дружеская шутка --- нарядить наивную мечту о самом большом числе в мантию
серьёзной науки, оставив саму математику (тетрация, логарифмическая шкала) честной.

И всё же «конечное» не значит «постижимое». Мальчики остановились на \(z\) лишь потому, что
кончился алфавит, а не числа; наибольшего числа нет --- любая вершина оказывается ступенью.
Взрослым эти числа «известны» --- но известно \emph{имя}, а не число: его не видел никто.
Посмейтесь над двумя мальчиками ещё раз. Уже не смешно, правда?

\paragraph{Приглашение к сотрудничеству.} Идеи, замечания и собственные восходящие ряды
присылайте на \href{mailto:k_mackl@mail.ru?subject=crypto\%20coma}{k\_mackl@mail.ru}
и \href{mailto:ink-shtil@mail.ru?subject=crypto\%20coma}{ink-shtil@mail.ru} с темой письма
\emph{crypto coma}.
